% =====================================================================
%  Lời mở đầu (TV1) — câu chuyện Hamming 1947 chuyển từ đầu Chương 1 sang
%  đây theo yêu cầu nhóm trưởng (11/11/2026). Chi tiết lịch sử: Bell
%  Telephone Laboratories 1947, máy tính rơ-le chạy cuối tuần không người
%  trực, máy bỏ dở công việc khi phát hiện lỗi; bài báo Hamming 1950 đã
%  kiểm chứng (IEEE Xplore, BSTJ 29(2), 147-160). Phần bố cục báo cáo bám
%  đúng tên 6 chương + Phụ lục A trong BTLDSTT.pdf.
% =====================================================================
\chuongphu{LỜI MỞ ĐẦU}

Một sáng thứ Hai năm 1947, tại phòng thí nghiệm Bell Telephone Laboratories
(Hoa Kỳ), nhà toán học Richard Hamming quay lại bên một máy tính rơ-le mà ông
đã giao chương trình chạy từ tối thứ Sáu. Dữ liệu vào là các bìa đục lỗ, và
suốt cuối tuần không có người trực máy. Ông chỉ thấy một khoảng trống: máy đã
phát hiện ra một lỗi, bỏ dở công việc rồi chuyển sang công việc kế tiếp, nên cả
cuối tuần tính toán không để lại kết quả nào. Chuyện ấy không chỉ xảy ra một
lần.

Điều làm Hamming bực bội không phải là máy có thể sai, mà là máy \emph{biết}
mình vừa sai nhưng chỉ có thể bỏ cuộc. Ông đặt ra một câu hỏi giản dị: nếu máy
đã nhận ra có một bit bị đảo, tại sao nó không xác định luôn bit đó nằm ở đâu
để tự sửa lại? Từ câu hỏi ấy ra đời bài báo \emph{Error Detecting and Error
Correcting Codes} \cite{hamming1950}, công bố năm 1950 trên tạp chí \emph{Bell
System Technical Journal}, và mã Hamming ra đời cùng với nó.

Cách sửa lỗi dễ hình dung nhất là \emph{mã lặp}: gửi mỗi bit ba lần rồi lấy
theo đa số. Cách này sửa được một lỗi nhưng rất tốn kém: tốc độ mã (tỉ số giữa
số bit dữ liệu và độ dài từ mã) chỉ là $R = 1/3$, tức mỗi bit dữ liệu phải
gánh thêm $2$ bit dư. Hamming chỉ ra rằng có thể làm tốt hơn nhiều: chỉ cần
$3$ bit kiểm tra để bảo vệ $4$ bit dữ liệu ($R = 4/7$), mà bên nhận vẫn tìm ra
và sửa được một bit sai bất kỳ.

Bí quyết nằm ở chỗ biến việc ``tìm bit sai'' thành một bài toán của đại số
tuyến tính trên trường $\GF$ chỉ có hai phần tử $0$ và $1$: các bit kiểm tra là
tổ hợp tuyến tính của các bit dữ liệu, và ``dấu vết'' mà một bit sai để lại
(gọi là \emph{hội chứng}) cho biết bit nào đã bị đảo. Đó cũng là sợi chỉ xuyên
suốt bài tập lớn này.

Bài tập lớn môn Đại số tuyến tính (MT1007) của nhóm 10 đi theo hướng ấy với đề
tài ``Mã Hamming (7,4): mô phỏng mã hóa, gây nhiễu, phát hiện và sửa lỗi cho
chuỗi ký tự''. Nhóm trình bày cơ sở lý thuyết, giải thích và chứng minh các
tính chất của ma trận tạo mã và ma trận kiểm tra, mở rộng sang mã Hamming
(8,4), rồi xây dựng phần mềm Python mô phỏng toàn bộ quá trình trên chuỗi ký
tự.

Báo cáo gồm sáu chương và một phụ lục. Chương~\ref{chap:1} trình bày cơ sở lý
thuyết của mã Hamming (7,4). Chương~\ref{chap:2} giải thích cơ sở lý thuyết của
các ma trận $G$, $H$ và mở rộng sang mã (8,4). Chương~\ref{chap:3} minh họa
bằng ví dụ tạo mã, kiểm tra lỗi và sửa lỗi. Chương~\ref{chap:4} trình bày thiết
kế và xây dựng phần mềm mô phỏng. Chương~\ref{chap:5} báo cáo kiểm thử và kết
quả thực nghiệm. Chương~\ref{chap:6} nêu kết luận và hướng phát triển. Phụ lục A
chứa mã nguồn chương trình Python.
